\clearpage\pagestyle{otMain}\setlength{\otBlocEspacement}{2.5mm}\fontsize{11}{15}\selectfont\corrigesChapitre{4}{Diviser et interpréter un partage}
\begin{blocLivre}{cor-4-1}\begin{corrige}{1}{Part ou nombre de parts ?}\demarche{Se demander ce que l’on cherche.}
1. Valeur d’une part : $96:8=\res{12}$ cartes par joueur. 2. Nombre de parts : $96:8=\res{12}$ paquets. Le calcul est le même, la question non.
\end{corrige}\end{blocLivre}
\begin{blocLivre}{cor-4-2}\begin{corrige}{2}{Division euclidienne}\demarche{Encadrer par des multiples du diviseur.}
$\res{487=23\times21+4}$, avec $4<23$. $\res{1\,000=7\times142+6}$, avec $6<7$.
\end{corrige}\end{blocLivre}
\begin{blocLivre}{cor-4-3}\begin{corrige}{3}{La bonne égalité}\demarche{Vérifier que le reste est plus petit que le diviseur.}
Les trois égalités sont justes, mais le reste doit être plus petit que 12 : seule $\res{158=12\times13+2}$ convient. Quotient 13, reste 2.
\end{corrige}\end{blocLivre}
\begin{blocLivre}{cor-4-4}\begin{corrige}{4}{Division décimale}\demarche{Diviser rang par rang en plaçant la virgule au passage des unités.}
$47{,}6:4=\res{11{,}9}$ ; $9{,}45:5=\res{1{,}89}$. Contrôles : $11{,}9\times4=47{,}6$ et $1{,}89\times5=9{,}45$.
\end{corrige}\end{blocLivre}
\begin{blocLivre}{cor-4-5}\begin{corrige}{5}{Après la virgule}\demarche{Transformer le reste en dixièmes, puis en centièmes.}
$23:4=\res{5{,}75}$. Pour $7:8$ : 70 dixièmes donnent 8 dixièmes, reste 6 ; 60 centièmes donnent 7 centièmes, reste 4 ; 40 millièmes donnent 5 millièmes, reste 0 : $\res{0{,}875}$.
\end{corrige}\end{blocLivre}
\begin{blocLivre}{cor-4-6}\begin{corrige}{6}{Une division posée}\demarche{Diviser les dizaines, puis les unités.}
Dans 36 dizaines, il y a 4 fois 8 : reste 4 dizaines. 45 unités contiennent 5 fois 8 : reste 5. Quotient $45$, reste $5$ : $\res{365=8\times45+5}$.
\label{fig-cor-4-6}\figureGeo[corrige]{div-6}
\end{corrige}\end{blocLivre}
\begin{blocLivre}{cor-4-7}\begin{corrige}{7}{Le reste dans l’histoire}\demarche{Interpréter le reste dans la situation.}
$50=12\times4+2$. Quatre minibus ne suffisent pas : il en faut $\res5$. Places libres : $5\times12-50=\res{10}$.
\end{corrige}\end{blocLivre}
\begin{blocLivre}{cor-4-8}\begin{corrige}{8}{Une division sans fin}\demarche{Repérer un reste qui se répète.}
Pour $10:3$, le reste 1 revient toujours : $3{,}333\ldots$ Pour $2:3$, le reste 2 revient : $0{,}666\ldots$ Les divisions ne s’arrêtent pas. Arrondis : $\res{3{,}33}$ et $\res{0{,}67}$.
\end{corrige}\end{blocLivre}
\begin{blocLivre}{cor-4-9}\begin{corrige}{9}{Partager au centime}\demarche{Diviser, puis contrôler par une multiplication.}
$25:4=\res{6{,}25}$ € : partage exact. $17:6=2{,}833\ldots$ : chacun reçoit $\res{2{,}83}$ € et il reste $17-6\times2{,}83=0{,}02$ €, soit 2 centimes.
\end{corrige}\end{blocLivre}
\begin{blocLivre}{cor-4-10}\begin{corrige}{10}{Retrouver le dividende}\demarche{Utiliser l’égalité dividende = diviseur × quotient + reste.}
Dividende : $15\times23+9=\res{354}$. Le reste ne peut pas valoir 15 : il doit être plus petit que le diviseur, sinon on pourrait faire un paquet de plus.
\end{corrige}\end{blocLivre}
\begin{blocLivre}{cor-4-11}\begin{corrige}{11}{Les boîtes d’œufs}\demarche{Quotient, puis interprétation du reste.}
$1\,000=6\times166+4$ : $\res{166}$ boîtes pleines et 4 œufs restants. Pour tout emballer, il faut une boîte de plus : $\res{167}$ boîtes.
\end{corrige}\end{blocLivre}
\begin{blocLivre}{cor-4-12}\begin{corrige}{12}{L’addition du restaurant}\demarche{Division décimale par un entier inférieur à 10.}
$94{,}5:7=\res{13{,}5}$, soit $13{,}50$ € chacun. Contrôle : $13{,}5\times7=94{,}5$.
\end{corrige}\end{blocLivre}
\begin{blocLivre}{cor-4-13}\begin{corrige}{13}{La randonnée}\demarche{Chercher la valeur d’une part.}
$38{,}4:3=\res{12{,}8}$ km par jour. $38{,}4:8=\res{4{,}8}$ km par étape.
\end{corrige}\end{blocLivre}
\begin{blocLivre}{cor-4-14}\begin{corrige}{14}{Quel jour sera-t-il ?}\demarche{Les semaines complètes ne changent pas le jour.}
$100=7\times14+2$ : 14 semaines complètes ramènent au lundi, puis 2 jours de plus. Ce sera un $\res{\text{mercredi}}$.
\end{corrige}\end{blocLivre}
\begin{blocLivre}{cor-4-15}\begin{corrige}{15}{Le plus grand reste}\demarche{Écrire le nombre sous la forme diviseur × quotient + reste.}
Le reste est plus petit que 24 : au plus $\res{23}$. Les nombres cherchés s’écrivent $24\times q+23$ : $24\times20+23=503$ dépasse 500, donc la réponse est $24\times19+23=\res{479}$.
\end{corrige}\end{blocLivre}
